% -------------------------------------------------------------------------
% ÜBERHOLT (31.08.2026). Maßgeblich ist die eingereichte main.tex. Zahlen,
% Zeilennummern und Angaben zum Overleaf-Stand in dieser Datei stammen aus
% der Zeit VOR der Crosswalk-Korrektur vom 29.08.2026 (35 statt 36 Stadt-
% teile, vollständiger Neulauf) und wurden bewusst NICHT nachgezogen: die
% Datei ist als datierter Entwurf Teil der Arbeitsdokumentation.
% -------------------------------------------------------------------------
% =========================================================================
% KAPITEL 5.3 -- KONSTRUKTION DER MERKMALE UND ZIELGROESSEN
% Fassung 1, 24.08.2026. Erster Schwerpunkt des Kapitels.
%
% GLEITOBJEKTE: docs/kapitel5-3_gleitobjekte.tex
%   Quelltext 2 (lst:quoten)    -> nach Absatz 3
%   Quelltext 3 (lst:krimindex) -> nach Absatz 5
%   Tabelle (tab:codebook)      -> nach Absatz 9
%
% NEUE ZITATE: Brantingham1997 (Absatz 4), Altman2006 (Absatz 10).
% Beide muessen vorher in literatur.bib stehen.
%
% ZEHN ABSAETZE STATT ACHT. Das Geruest sah acht vor und begann bei den
% Quoten. Dazugekommen sind der Einstieg (Bauplan wie 5.1 und 5.2) und die
% Aufteilung des Kriminalitaetsindex auf zwei Absaetze -- Definition und
% Quelle im einen, die beiden Begruendungen im anderen. Der Absatz waere
% sonst der laengste des Kapitels geworden.
%
% WAS BEWUSST NICHT DRINSTEHT
%   - Klassenverteilung 79,0 / 16,3 / 3,1 / 1,5 Prozent -> steht in 4.2
%     und 7.2. Verteilungskennzahlen gehoeren nicht nach Kapitel 5.
%   - Mittelwert, Median, Maximum der Mengenzielgroessen -> Tabelle
%     tab:kennzahlen in 4.2. Hier nur ein Verweis.
%   - Dispersionsindex 62,8 -> 4.2, wird in 5.4 fuer die Baseline zitiert.
%   - Die frueher hier vorgesehene Zielgroessentabelle mit Kennwerten
%     entfaellt damit; tab:codebook fuehrt Wertebereiche, keine Statistik.
%
% DIE VERWORFENEN AGGREGATIONSEBENEN aus 5.1 Fassung 1 sind untergebracht:
% Census Tract mal Jahr in Absatz 1, die Einzeleinsatzebene in Absatz 10 --
% dort mit gemessener Zahl statt als Behauptung.
% =========================================================================


% ---------------------------- AB HIER KOPIEREN ---------------------------

\subsection{Konstruktion der Merkmale und Zielgrößen}
\label{sec:merkmale}

Die Einsatztabelle aus Abschnitt~\ref{sec:integration} führt 719.989 Zeilen,
eine je Ereignis. Auf dieser Ebene stehen Merkmale und Zielgröße nicht auf
derselben Beobachtung: Die Strukturmerkmale beschreiben einen Stadtteil, die
Zielgröße soll dessen Einsatzlast beschreiben. Die Verdichtung auf Stadtteil
und Monat stellt beides in dieselbe Zeile und bildet dabei die Größen, mit
denen gerechnet wird. Eine Verdichtung auf Census Tract und Jahr scheidet aus,
weil die Einsatzdaten keine Tract-Zuordnung tragen
(Abschnitt~\ref{sec:datengrundlage}) und die Saisonalität dabei verloren ginge.
Tabelle~\ref{tab:codebook} führt die entstehenden Größen mit Skalenniveau,
Einheit und Wertebereich; die folgenden Absätze begründen ihre Bildung. Am Ende
stehen zwei Dateien: 4.620 Zeilen für die Menge der Einsatzlast und 4.619 für
ihre Zusammensetzung.

Verdichtet wird auf ein vollständiges Raster aus 35 Stadtteilen und 132
Monaten. Vollständig heißt, dass auch ein Monat ohne Einsatz eine Zeile erhält,
und zwar mit dem Wert null statt gar nicht. Ohne dieses Raster fehlte ein
ruhiger Monat stillschweigend, und jedes rückwärtsgerichtete Merkmal griffe um
eine Position zu weit zurück. Fehlende Merkmalswerte werden ausschließlich
vorwärts fortgeschrieben; ein Rückwärtsfüllen schlösse eine Lücke mit einem
später erhobenen Wert und spielte damit genau die Information ein, die zum
Prognosezeitpunkt nicht vorlag (Abschnitt~\ref{sec:auswahl-bereinigung}). Die
4.620 Zeilen sind damit 35 mal 132, lückenlos. Dem Strukturdatensatz fehlt der
eine Monat ohne Einsatz, für den sich keine Zusammensetzung bilden lässt.

Sechs der zehn Strukturmerkmale sind Anteilswerte, gebildet aus je zwei
Zählgrößen derselben Erhebung. Ein Nenner von null ergibt dabei einen fehlenden
Wert und nicht null, damit eine fehlende Bezugsgröße sichtbar bleibt, statt als
Anteil von null in die Modellierung zu gehen (Quelltext~\ref{lst:quoten}). Ein
Nenner ist begründungspflichtig: Der Altbauanteil bezieht sich auf die Parzellen
mit bekanntem Baujahr und nicht auf alle Parzellen. Ein Teil des Verzeichnisses
trägt kein Baujahr (Abschnitt~\ref{sec:eda}); mit allen Parzellen im Nenner
fiele der Altbauanteil systematisch zu niedrig aus, und zwar dort am stärksten,
wo die Erfassung am lückenhaftesten ist.

Das kriminalitätsbezogene Merkmal ist kein Rohwert, sondern ein
Standortquotient: die Deliktrate eines Stadtteils, gemessen an der Deliktrate
der Gesamtstadt im selben Monat.

\begin{equation*}
  \mathrm{Rate}(i,t) =
  \frac{\mathrm{Delikte}(i,\ \text{Fenster endend in } t-1)}{\mathrm{Einwohner}(i)}
  \qquad
  \mathrm{Index}(i,t) = \frac{\mathrm{Rate}(i,t)}{\mathrm{Rate}(\mathrm{Stadt},t)}
\end{equation*}

Ein Wert von eins bedeutet eine Belastung wie im Stadtdurchschnitt desselben
Monats. Die Form des Standortquotienten stammt aus der Regionalwissenschaft und
wird von Brantingham und Brantingham auf Kriminalität
übertragen.\footcite[268-269]{Brantingham1997} Deren Fassung verwendet
allerdings die Gesamtzahl der Delikte als Bezugsgröße und misst damit, worauf
ein Gebiet relativ spezialisiert ist. Hier ist die Bezugsgröße die
Wohnbevölkerung; gemessen wird die relative Belastung je Einwohner und nicht
die Zusammensetzung der Delikte.

Zwei Eigenschaften dieser Konstruktion sind Entscheidungen. Die relative Form
ist eine Antwort auf den Quellenwechsel aus Abschnitt~\ref{sec:integration}: Ein
Wechsel des Erfassungssystems verschiebt das stadtweite Niveau annähernd
multiplikativ, wirkt damit auf Zähler und Nenner gleichermaßen und kürzt sich
weitgehend heraus. Was sich nicht herauskürzt, ist eine Verschiebung in der
Zusammensetzung der erfassten Delikte, die einzelne Stadtteile stärker trifft
als andere; diese Verzerrung bleibt bestehen. Die zweite Eigenschaft ist das
Fenster: Es endet strikt im Vormonat, weil eine im selben Monat gemessene
Kriminalität die Einsätze dieses Monats erklären würde, statt sie
vorherzusagen (Quelltext~\ref{lst:krimindex}). Der Index geht logarithmiert
ein, sodass der Stadtdurchschnitt bei null liegt und gleich große Abweichungen
nach oben und unten denselben Abstand haben.

Die Wohnbevölkerung geht als eigenes Merkmal ein, ebenfalls logarithmiert. Sie
wirkt auf die beiden Mengenzielgrößen gegenläufig: Ein größerer Stadtteil hat
mehr Einsätze, aber nicht zwangsläufig mehr Einsätze je Einwohner. Ohne diese
Größenkontrolle bildete ein Modell auf der absoluten Zielgröße vor allem die
Einwohnerzahl ab und nicht die Struktur, nach der gefragt ist. Im Poisson-GLM
des Abschnitts~\ref{sec:rahmen-baselines} tritt sie zusätzlich als Offset auf.

Der Kalendermonat wird nicht als Zahl von eins bis zwölf geführt, sondern als
Sinus- und Kosinusterm. Als Zahl hätten Dezember und Januar den Abstand elf,
obwohl sie aufeinanderfolgen; die beiden Terme bilden den Jahreszyklus dagegen
geschlossen ab. Sie sind die einzigen Merkmale, die ausschließlich innerhalb
der Stadtteile variieren (Abschnitt~\ref{sec:eda}).

Der Datensatz führt zusätzlich drei Verlaufsmerkmale: den Wert des Vormonats,
den des Vorjahresmonats und das Mittel der drei Vormonate. Sie sind strikt
rückwärtsgerichtet und werden je Stadtteil gebildet, nie über Stadtteilgrenzen
hinweg; beim gleitenden Mittel wird zuerst verschoben und dann gemittelt, sodass
der Wert für einen Monat diesen Monat selbst nicht enthält. Damit sie schon für
den ersten Analysemonat definiert sind, beginnt die Aggregation zwölf Monate
früher; diese Vorlaufmonate gehen ausschließlich über die Verschiebung ein und
bilden keine eigenen Zeilen. In die Modelle gehen die Verlaufsmerkmale dennoch
nicht ein. Unter der Aufteilung des Abschnitts~\ref{sec:rahmen-baselines} wird
ein Stadtteil vollständig zurückgehalten; sein Vormonatswert wäre seine eigene
Vergangenheit, und das Modell beantwortete dann die Frage, wie sich ein
bekannter Stadtteil fortschreibt, statt der Frage, was sich über einen
unbekannten sagen lässt. Sie verbleiben im Datensatz, weil die Beschreibung in
Kapitel~\ref{sec:datengrundlage} sie benötigt.

Drei Zielgrößen stehen nebeneinander (Abschnitt~\ref{sec:zielgroessen}): die
Zahl der Einsätze je Stadtteilmonat, dieselbe Zahl auf je 1.000 Einwohner
normiert und die überwiegende Einsatzart. Die vier Anteile der NFIRS-Gruppen an
der Einsatzlast eines Monats sind dabei Zwischengröße und keine eigenen
Zielgrößen: Aus ihnen entsteht die überwiegende Einsatzart als der größte der
vier. Sie werden weder modelliert noch berichtet. Lage und Streuung der beiden
Mengenzielgrößen stehen in Abschnitt~\ref{sec:eda}.

Dass die Einsatzart auf Stadtteilebene und nicht je Einsatz bestimmt wird, ist
die folgenreichste Festlegung dieses Abschnitts. Innerhalb eines Stadtteilmonats
tragen alle Einsätze dieselben Strukturmerkmale; die 350.481 Einsätze des
Analysezeitraums verteilen sich deshalb auf nur 4.619 verschiedene
Merkmalsprofile --- genau so viele, wie es Stadtteilmonate mit mindestens einem
Einsatz gibt. Ein Modell, das jedem Profil die dort häufigste Einsatzart zuwiese
und auf dieser Ebene nicht mehr zu verbessern wäre, träfe in 49,9\,\% der Fälle
zu; die häufigste Klasse ohne jedes Merkmal zu nennen trifft in 48,2\,\%. Der
gesamte Spielraum beträgt 1,7 Prozentpunkte. Auf Stadtteilebene ist die Frage
dagegen beantwortbar. Die Zielgröße bleibt dabei eine echte Klasse und kein an
einem Schwellwert geteilter stetiger Anteil; damit entfällt die
Begründungslast einer solchen Teilung, die Trennschärfe kostet und deren Wirkung
Altman und Royston mit dem Verlust eines Drittels der Beobachtungen
beziffern.\footcite{Altman2006}

% ---------------------------- BIS HIER KOPIEREN --------------------------
